\section{温度关联分析与因果识别边界}\label{sec:temperature-association}

\subsection{温度干预问题与假设图}
对给定的完整药物形式和溶剂组成，因果问题是按同一实验协议将体系从温度 $t_0$ 调至 $t_1$ 并达到相应平衡时，对数溶解度会改变多少。用 $D$ 和 $S$ 表示观测到的药物与溶剂，$T$ 和 $Y$ 表示记录温度和对数溶解度，$U$ 表示未记录的实验条件。以下有向关系给出识别所依赖的假设\cite{pearl2009}：
\begin{equation}
(D,S)\longrightarrow T,\qquad (D,S)\longrightarrow Y,\qquad
U\longrightarrow T,\qquad U\longrightarrow Y,\qquad T\longrightarrow Y.
\end{equation}
药物与溶剂可能同时影响温度选择和溶解度，实验来源或协议也可能同时影响二者。$T\to Y$ 表示拟考察的因果路径。温度字段沿用 CSV 原始数值尺度，以下斜率以每一原始温度单位对应的对数溶解度变化表示。

\subsection{总体关联与组合内关联}
在开发集中，以规范化的完整药物与完整溶剂结构定义组合，保留每个组合内至少两个不同温度的记录。得到 157 个组合、73 个药物母体和 1118 条记录。各组合内温度跨度为 5.00 至 65.15 个原始单位，中位数为 27.10。在相同的 1118 条记录上，先按逐行线性回归估计总体斜率，再以组合固定效应估计共同的组内斜率。令 $c$ 表示组合，$i$ 表示组合内记录，则
\begin{align}
Y_{ic}&=\alpha_c+\beta T_{ic}+\varepsilon_{ic},\\
\hat\beta&=\frac{\sum_c\sum_i(T_{ic}-\bar T_c)(Y_{ic}-\bar Y_c)}
{\sum_c\sum_i(T_{ic}-\bar T_c)^2}.
\end{align}
组合内去均值可消去 $\alpha_c$，使斜率由同一组合内的温度和溶解度共同变化决定。若各组合的斜率不同，共同斜率相当于按组合内温度离差平方和加权，温区较宽或记录较多的组合可占较大权重。

\subsection{估计结果与不确定性}
同一子集上，总体斜率点估计为负，组内点估计为正（表~\ref{tab:temperature-association}）。总体区间覆盖零。在相同药物与溶剂的记录之间比较时，较高温度与较高溶解度相联系。按母体进行 2000 次有放回重采样，各次均保留被抽中母体的全部组合和记录，得到组内斜率的 95\% 百分位区间 $[0.01236,0.01902]$。逐一删去母体后的斜率范围为 $[0.01424,0.01599]$，各次估计保持正向。

\begin{table}[htbp]
\centering\small
\caption{相同 1118 条记录上的总体和组合内斜率。单位为每一 CSV 原始温度单位对应的对数溶解度变化，95\% 区间采用母体聚类稳健标准误。}
\label{tab:temperature-association}
\begin{tabular}{lrr}
\toprule
估计口径 & 斜率 & 95\% 区间\\
\midrule
逐行总体回归 & $-0.00746$ & $[-0.02614,\phantom{-}0.01123]$\\
组合内固定效应 & $\phantom{-}0.01538$ & $[\phantom{-}0.01198,\phantom{-}0.01878]$\\
\bottomrule
\end{tabular}
\end{table}

\subsection{未测混杂与识别条件}
固定效应控制组合内稳定的差异，实验来源、协议等若随温度记录变化，仍可形成温度与溶解度的共同影响。pH 和晶型的角色取决于测量顺序及干预条件。若它们随温度变化，固定或调整它们会改变所对应的因果问题。在去组均值线性模型中，进一步衡量未测变量对斜率点估计的影响\cite{sensitivity2020}。记 $r_T$ 为该变量与温度在控制组合后的偏 $R^2$，$r_Y$ 为该变量与溶解度在控制温度与组合后的偏 $R^2$。在设定 $r_T=r_Y$ 且偏倚方向使观测斜率增大的情形下，两者约需同为 0.54 才能使调整后的斜率降为零。该数值用普通最小二乘的遗漏变量偏倚恒等式计算，刻画点估计对假设混杂强度的敏感性。

要将组内关联解释为干预效应，需核实温标，在可比协议下明确温度与平衡状态，并在有共同温度支持的组合内测量或控制先于干预的混杂因素。当前估计描述实测组合内的关联，后续可通过受控温度实验检验这一关系。
